% Заключительные слайды презентации

\section{Выводы}

\begin{frame}{Основные результаты}
    \begin{itemize}
        \item Предложена новая формальная модель кибератак — \textbf{атрибутивный метаграф}, объединяющая семантику событий и тактики MITRE ATT\&CK.
        \item Разработан и реализован метод обнаружения, превосходящий современные аналоги: \textbf{F1-score = 0.87} (+23\% относительно SOTA).
        \item Подтверждена применимость на \textbf{реальных данных}: APT29, финансовые атаки, корпоративный трафик MIREA.
        \item Результаты опубликованы в журнале «Проблемы информационной безопасности. Компьютерные системы».
    \end{itemize}
\end{frame}

\begin{frame}{Практическая значимость}
    \begin{itemize}
        \item Метод внедрён в пилотный модуль системы управления инцидентами безопасности (СУИБ) одного из российских банков.
        \item Архитектура решения поддерживает потоковую обработку и масштабирование на кластере.
        \item Открытый код алгоритма доступен для научного сообщества (по запросу).
    \end{itemize}
\end{frame}

\begin{frame}{Публикации по теме диссертации}
    \begin{itemize}
        \item Латыпов И.Т., Еремеев М.А. Обнаружение APT-атак на основе атрибутивных метаграфов // \textit{Проблемы информационной безопасности. Компьютерные системы}. — 2025. — №2. — С. 45–58. \hfill \textcolor{blue}{(ВАК)}
        \item Латыпов И.Т., Еремеев М.А. Подход к обнаружению целевых атак на начальном этапе с использованием модели компьютерной атаки на базе атрибутивных метаграфов // \textit{Кибербезопасность: технические и правовые аспекты защиты информации}. — М., 2024. — С. 308–312. \hfill \textcolor{blue}{(конф.)}
        \item Латыпов И.Т., Криулин А.А. Анализ действий APT группировок при проведении целевых атак с использованием модели выявления вредоносной активности на базе атрибутивных метаграфов // \textit{Кибербезопасность: технические и правовые аспекты защиты информации}. — М., 2023. — С. 48–51. \hfill \textcolor{blue}{(конф.)}
        \item Смирнов С.И., Изергин Д.А., Латыпов И.Т. Программный комплекс автоматизации процесса извлечения признаков исполняемых файлов ОС Android... — Свидетельство о регистрации ПО № RU 2021612358. — 2021. \hfill \textcolor{blue}{(ПО)}
    \end{itemize}
\end{frame}

\begin{frame}{Благодарю за внимание!}
  \centering
  \vspace{2cm}
  \Large Вопросы?
\end{frame}

\begin{comment}
\begin{frame}
    \frametitle{Научная новизна}
    \begin{itemize}
        \item Впервые реализован \dots
        \item Разработана программа \dots
        \item Впервые проведён анализ \dots
        \item Предложена схема \dots
    \end{itemize}
\end{frame}
\note{
    Проговаривается вслух научная новизна
}

\begin{frame}
    \frametitle{Научная и практическая значимость}
    \begin{itemize}
        \item Получены выражения для \dots.
        \item Определены условия \dots.
        \item Разработаны устройства \dots.
    \end{itemize}
\end{frame}
\note{
    Проговариваются вслух научная и практическая значимость
}

\begin{frame}
    \frametitle{Свидетельство о регистрации программы}
    \begin{figure}[h]
        \centering
        \includegraphics[height=0.7\textheight]{registration}
    \end{figure}
\end{frame}
\note{
    Получено свидетельство о регистрации разработанной программы \textsc{Hello~world™}.
}

\begin{frame}
    \frametitle{Акт о внедрении}
    \begin{figure}[h]
        \centering
        \fbox{
            \begin{minipage}[t]{0.4\linewidth}
                \includegraphics[width=\linewidth]{implementation}
            \end{minipage}
        }
    \end{figure}
\end{frame}
\note{
    Получен акт о внедрении.
}

\begin{frame} % публикации на одной странице
% \begin{frame}[t,allowframebreaks] % публикации на нескольких страницах
    \frametitle{Основные публикации}
    \nocite{vakbib1}%
    \nocite{vakbib2}%
    %
    %% authorwos
    \nocite{wosbib1}%
    %
    %% authorscopus
    \nocite{scbib1}%
    %
    %% authorconf
    \nocite{confbib1}%
    \nocite{confbib2}%
    %
    %% authorother
    \nocite{bib1}%
    \nocite{bib2}%
    \ifnumequal{\value{bibliosel}}{0}{
        \insertbiblioauthor
    }{
        \printbibliography%
    }
\end{frame}
\note{
    Результаты работы опубликованы в N печатных изданиях,
    в~т.\:ч. M реферируемых изданиях.
}

\begin{frame}
    \frametitle{Участие в конференциях}
    \begin{itemize}
        \item Научная сессия МГУ, Москва 2013--2015;
        \item \rom{24} Russian Conference (RuC 2014), Obninsk, Russia, 2014
        \item \rom{7} International Conference (IAC 16), Busan, Korea,
              2016;
        \item \rom{28} Other Conference (AC 16), East Lansing, MI USA, 2016;
        \item \dots
    \end{itemize}
\end{frame}
\note{
    Работа была представлена на ряде конференций.
}

\begin{frame}[plain, noframenumbering] % последний слайд без оформления
    \begin{center}
        \Huge
        Спасибо за внимание!
    \end{center}
\end{frame}
\end{comment}